% 第十一章 附录

\chapter{附录}

\section{附录A：120+细分板块景气度评分明细}

本附录列示本报告覆盖的全部 120+ 细分板块景气度评分结果。评分基于"业绩验证---产业映射---政策加持---估值持仓"四维度框架，采用百分制。板块分类基于"景气度---拥挤度"二维矩阵，分为明星型、潜力型、博弈型、陷阱型四类。

\vspace{0.3em}

为便于查阅，按一级行业分组列示，组内按景气度评分降序排列。完整评分模型及权重设置详见附录 E。

\begin{exhibitbox}[表 A-1：TMT 行业细分板块景气度评分明细]
\centering
\scriptsize
\renewcommand{\arraystretch}{1.1}
\begin{tabular*}{\textwidth}{@{\extracolsep{\fill}}l r r r c l@{}}
\toprule
\textbf{细分板块} & \textbf{景气度评分} & \textbf{业绩增速 (2026Q1)} & \textbf{拥挤度评分} & \textbf{景气趋势} & \textbf{板块分类} \\
\midrule
AI 算力芯片 & 93.5 & +62.3\% & 88 & $\nearrow$ 加速上行 & 明星型 \\
光模块 & 92.1 & +58.7\% & 92 & $\nearrow$ 加速上行 & 明星型 \\
功率半导体 & 89.4 & +35.2\% & 58 & $\nearrow$ 持续上行 & 潜力型 \\
半导体设备 & 87.8 & +32.6\% & 75 & $\nearrow$ 持续上行 & 潜力型 \\
半导体材料 & 83.2 & +25.8\% & 62 & $\nearrow$ 持续上行 & 潜力型 \\
存储芯片 & 78.5 & +22.4\% & 70 & $\rightarrow$ 高位平稳 & 明星型 \\
模拟芯片 & 72.3 & +15.6\% & 55 & $\nearrow$ 缓慢上行 & 潜力型 \\
数字芯片设计 & 70.8 & +18.9\% & 68 & $\rightarrow$ 平稳 & 博弈型 \\
封测 & 68.5 & +12.3\% & 48 & $\nearrow$ 缓慢上行 & 潜力型 \\
分立器件 & 65.2 & +8.7\% & 42 & $\rightarrow$ 平稳 & 潜力型 \\
消费电子芯片 & 58.6 & -3.2\% & 38 & $\searrow$ 下行 & 陷阱型 \\
\midrule
\multicolumn{6}{l}{\kaishu{（半导体子行业，共 11 个细分板块）}} \\
\midrule
智驾零部件 & 84.6 & +33.4\% & 65 & $\nearrow$ 加速上行 & 潜力型 \\
智能座舱 & 76.8 & +24.5\% & 60 & $\nearrow$ 持续上行 & 潜力型 \\
车载半导体 & 75.3 & +28.7\% & 58 & $\nearrow$ 持续上行 & 潜力型 \\
汽车电子 & 72.1 & +19.8\% & 52 & $\nearrow$ 持续上行 & 潜力型 \\
汽车线束 & 62.5 & +10.2\% & 35 & $\rightarrow$ 平稳增长 & 潜力型 \\
\midrule
\multicolumn{6}{l}{\kaishu{（汽车电子子行业，共 5 个细分板块）}} \\
\midrule
AI 大模型 & 82.7 & +45.6\% & 85 & $\nearrow$ 加速上行 & 明星型 \\
云计算 & 78.2 & +28.9\% & 72 & $\nearrow$ 持续上行 & 明星型 \\
工业软件 & 74.5 & +21.3\% & 58 & $\nearrow$ 持续上行 & 潜力型 \\
信创 & 70.5 & +18.5\% & 65 & $\rightarrow$ 平稳 & 博弈型 \\
网络安全 & 68.3 & +15.2\% & 48 & $\nearrow$ 缓慢上行 & 潜力型 \\
IT 服务 & 56.8 & +5.3\% & 32 & $\rightarrow$ 低位平稳 & 陷阱型 \\
\midrule
\multicolumn{6}{l}{\kaishu{（计算机子行业，共 6 个细分板块）}} \\
\midrule
光纤光缆 & 72.1 & +15.4\% & 45 & $\rightarrow$ 平稳增长 & 潜力型 \\
通信设备 & 68.7 & +12.8\% & 42 & $\rightarrow$ 平稳 & 潜力型 \\
运营商 & 65.4 & +8.9\% & 38 & $\rightarrow$ 平稳增长 & 潜力型 \\
6G & 52.3 & +3.5\% & 78 & $\rightarrow$ 主题驱动 & 博弈型 \\
卫星互联网 & 48.6 & -2.1\% & 75 & $\rightarrow$ 主题驱动 & 博弈型 \\
\midrule
\multicolumn{6}{l}{\kaishu{（通信子行业，共 5 个细分板块）}} \\
\bottomrule
\end{tabular*}
\vspace{0.5em}
\sourcenote{Wind，行业协会，本研究团队测算}
\end{exhibitbox}

\begin{exhibitbox}[表 A-2：新能源与电力设备行业细分板块景气度评分明细]
\centering
\scriptsize
\renewcommand{\arraystretch}{1.1}
\begin{tabular*}{\textwidth}{@{\extracolsep{\fill}}l r r r c l@{}}
\toprule
\textbf{细分板块} & \textbf{景气度评分} & \textbf{业绩增速 (2026Q1)} & \textbf{拥挤度评分} & \textbf{景气趋势} & \textbf{板块分类} \\
\midrule
储能电池 & 86.7 & +41.5\% & 82 & $\nearrow$ 持续上行 & 明星型 \\
动力电池 & 77.3 & +12.8\% & 38 & $\nearrow$ 拐点向上 & 潜力型 \\
\midrule
\multicolumn{6}{l}{\kaishu{（储能与电池子行业，共 2 个细分板块）}} \\
\midrule
光伏逆变器 & 82.5 & +31.2\% & 72 & $\rightarrow$ 高位平稳 & 明星型 \\
光伏组件 & 68.5 & +5.6\% & 35 & $\rightarrow$ 底部回升 & 潜力型 \\
光伏硅料 & 62.8 & -15.3\% & 28 & $\searrow$ 底部震荡 & 陷阱型 \\
\midrule
\multicolumn{6}{l}{\kaishu{（光伏子行业，共 3 个细分板块）}} \\
\midrule
电力运营 & 78.8 & +15.6\% & 28 & $\nearrow$ 持续上行 & 潜力型 \\
新能源发电 & 73.5 & +18.9\% & 30 & $\nearrow$ 持续上行 & 潜力型 \\
火电 & 65.7 & +9.2\% & 25 & $\nearrow$ 缓慢上行 & 潜力型 \\
水电 & 62.4 & +6.8\% & 22 & $\rightarrow$ 平稳 & 陷阱型 \\
\midrule
\multicolumn{6}{l}{\kaishu{（电力运营子行业，共 4 个细分板块）}} \\
\midrule
特高压 & 74.8 & +21.7\% & 42 & $\rightarrow$ 平稳增长 & 潜力型 \\
充电桩 & 72.6 & +24.3\% & 48 & $\nearrow$ 持续上行 & 潜力型 \\
智能电网 & 70.3 & +16.5\% & 38 & $\nearrow$ 持续上行 & 潜力型 \\
\midrule
\multicolumn{6}{l}{\kaishu{（电网设备子行业，共 3 个细分板块）}} \\
\midrule
风电零部件 & 67.8 & +12.4\% & 38 & $\nearrow$ 持续上行 & 潜力型 \\
风电整机 & 65.2 & +8.9\% & 32 & $\nearrow$ 缓慢上行 & 潜力型 \\
\bottomrule
\end{tabular*}
\vspace{0.5em}
\sourcenote{Wind，国家能源局，中电联，本研究团队测算}
\end{exhibitbox}

\begin{exhibitbox}[表 A-3：医药与高端制造行业细分板块景气度评分明细]
\centering
\scriptsize
\renewcommand{\arraystretch}{1.1}
\begin{tabular*}{\textwidth}{@{\extracolsep{\fill}}l r r r c l@{}}
\toprule
\textbf{细分板块} & \textbf{景气度评分} & \textbf{业绩增速 (2026Q1)} & \textbf{拥挤度评分} & \textbf{景气趋势} & \textbf{板块分类} \\
\midrule
创新药 & 85.3 & +28.9\% & 42 & $\nearrow$ 拐点向上 & 潜力型 \\
医疗器械 & 81.9 & +18.7\% & 35 & $\nearrow$ 拐点向上 & 潜力型 \\
CXO & 72.4 & +15.3\% & 48 & $\rightarrow$ 平稳 & 潜力型 \\
医疗服务 & 68.5 & +12.4\% & 42 & $\nearrow$ 缓慢上行 & 潜力型 \\
生物制品 & 67.2 & +10.5\% & 40 & $\nearrow$ 缓慢上行 & 潜力型 \\
中药 & 62.3 & +7.8\% & 35 & $\rightarrow$ 平稳 & 陷阱型 \\
原料药 & 58.9 & +2.1\% & 30 & $\searrow$ 下行 & 陷阱型 \\
医药商业 & 55.7 & +4.2\% & 28 & $\rightarrow$ 平稳 & 陷阱型 \\
\midrule
\multicolumn{6}{l}{\kaishu{（医药生物子行业，共 8 个细分板块）}} \\
\midrule
工业机器人 & 79.4 & +24.5\% & 55 & $\nearrow$ 持续上行 & 潜力型 \\
高端机床 & 75.2 & +19.3\% & 48 & $\nearrow$ 持续上行 & 潜力型 \\
检测服务 & 71.3 & +12.6\% & 40 & $\nearrow$ 持续上行 & 潜力型 \\
激光设备 & 68.7 & +15.8\% & 42 & $\nearrow$ 缓慢上行 & 潜力型 \\
轨交设备 & 62.5 & +8.4\% & 32 & $\rightarrow$ 平稳 & 潜力型 \\
工程机械 & 55.3 & -5.2\% & 28 & $\rightarrow$ 底部震荡 & 陷阱型 \\
\midrule
\multicolumn{6}{l}{\kaishu{（机械设备子行业，共 6 个细分板块）}} \\
\midrule
军工电子 & 80.6 & +22.3\% & 62 & $\rightarrow$ 高位平稳 & 明星型 \\
航空航天 & 76.5 & +16.9\% & 58 & $\rightarrow$ 平稳增长 & 潜力型 \\
军工新材料 & 74.2 & +18.5\% & 52 & $\nearrow$ 持续上行 & 潜力型 \\
地面装备 & 65.8 & +9.6\% & 38 & $\rightarrow$ 平稳增长 & 潜力型 \\
\midrule
\multicolumn{6}{l}{\kaishu{（国防军工子行业，共 4 个细分板块）}} \\
\bottomrule
\end{tabular*}
\vspace{0.5em}
\sourcenote{Wind，药监局，行业协会，本研究团队测算}
\end{exhibitbox}

\begin{exhibitbox}[表 A-4：其他行业细分板块景气度评分明细（节选）]
\centering
\scriptsize
\renewcommand{\arraystretch}{1.1}
\begin{tabular*}{\textwidth}{@{\extracolsep{\fill}}l r r r c l@{}}
\toprule
\textbf{细分板块} & \textbf{景气度评分} & \textbf{业绩增速 (2026Q1)} & \textbf{拥挤度评分} & \textbf{景气趋势} & \textbf{板块分类} \\
\midrule
合成生物学 & 73.5 & +18.2\% & 55 & $\nearrow$ 加速上行 & 潜力型 \\
人形机器人 & 68.9 & +8.5\% & 75 & $\nearrow$ 主题驱动 & 博弈型 \\
新型显示 & 65.2 & +12.3\% & 48 & $\nearrow$ 缓慢上行 & 潜力型 \\
低空经济 & 45.7 & -5.8\% & 82 & $\rightarrow$ 主题驱动 & 博弈型 \\
量子计算 & 42.3 & -10.2\% & 70 & $\rightarrow$ 主题驱动 & 博弈型 \\
\midrule
\multicolumn{6}{l}{\kaishu{（新兴主题板块，共 5 个）}} \\
\midrule
白酒 & 63.8 & +8.5\% & 45 & $\rightarrow$ 平稳 & 陷阱型 \\
啤酒 & 58.5 & +3.2\% & 38 & $\rightarrow$ 平稳 & 陷阱型 \\
乳制品 & 55.6 & +2.1\% & 32 & $\rightarrow$ 平稳 & 陷阱型 \\
调味品 & 52.3 & -2.5\% & 30 & $\searrow$ 下行 & 陷阱型 \\
\midrule
\multicolumn{6}{l}{\kaishu{（食品饮料子行业，共 4 个细分板块）}} \\
\midrule
保险 & 58.9 & +5.2\% & 32 & $\nearrow$ 缓慢回升 & 陷阱型 \\
银行 & 55.7 & +3.8\% & 28 & $\rightarrow$ 平稳 & 陷阱型 \\
证券 & 52.3 & -8.5\% & 42 & $\rightarrow$ 低位震荡 & 陷阱型 \\
\midrule
\multicolumn{6}{l}{\kaishu{（金融子行业，共 3 个细分板块）}} \\
\midrule
\multicolumn{6}{l}{\kaishu{注：完整 120+ 细分板块评分详见在线数据平台，此处列示重点及代表性板块。}} \\
\bottomrule
\end{tabular*}
\vspace{0.5em}
\sourcenote{Wind，本研究团队测算}
\end{exhibitbox}

\vspace{0.3em}

\textbf{评分说明}：景气度评分范围 0--100 分，基于业绩维度（40\%）、产业维度（30\%）、政策维度（15\%）、资金维度（15\%）加权计算。拥挤度评分范围 0--100 分，基于公募持仓拥挤度（40\%）、交易拥挤度（30\%）、估值分位拥挤度（30\%）加权计算。具体方法详见附录 E。

\section{附录B：核心标的财务预测与估值表}

本附录列示本报告覆盖的七大重点板块核心标的财务预测与估值数据。所有盈利预测基于 Wind 一致预期及本研究团队调整，市场数据截至 2026 年 6 月 20 日。

\begin{exhibitbox}[表 B-1：功率半导体板块核心标的财务预测与估值]
\centering
\scriptsize
\renewcommand{\arraystretch}{1.1}
\begin{tabular*}{\textwidth}{@{\extracolsep{\fill}}l r r r r r r c@{}}
\toprule
\textbf{公司} & \textbf{市值} & \textbf{2026E 营收} & \textbf{2026E 净利} & \textbf{2026E PE} & \textbf{2027E PE} & \textbf{2026E PEG} & \textbf{评级} \\
& (亿元) & (亿元/增速) & (亿元/增速) & (倍) & (倍) & (倍) & \\
\midrule
斯达半导 & 856 & 58.5 / 38\% & 16.2 / 42\% & 52.8 & 36.5 & 1.26 & \textcolor{riskgreen}{超配} \\
时代电气 & 782 & 245.0 / 22\% & 38.5 / 35\% & 20.3 & 15.6 & 0.58 & \textcolor{riskgreen}{超配} \\
士兰微 & 625 & 142.0 / 28\% & 15.8 / 65\% & 39.6 & 26.8 & 0.61 & \textcolor{riskgreen}{超配} \\
闻泰科技 & 598 & 765.0 / 15\% & 28.5 / 32\% & 21.0 & 16.2 & 0.66 & \textcolor{navy}{标配} \\
扬杰科技 & 286 & 58.2 / 25\% & 9.5 / 38\% & 30.1 & 22.4 & 0.79 & \textcolor{navy}{标配} \\
新洁能 & 245 & 28.5 / 32\% & 6.8 / 45\% & 36.0 & 24.8 & 0.80 & \textcolor{navy}{标配} \\
华润微 & 512 & 118.0 / 24\% & 18.2 / 40\% & 28.1 & 20.8 & 0.70 & \textcolor{navy}{标配} \\
东微半导 & 168 & 18.6 / 42\% & 4.8 / 55\% & 35.0 & 22.6 & 0.64 & \textcolor{riskgreen}{超配} \\
\bottomrule
\end{tabular*}
\vspace{0.5em}
\sourcenote{Wind，卖方一致预期，本研究团队预测；市场数据截至 2026 年 6 月 20 日}
\end{exhibitbox}

\begin{exhibitbox}[表 B-2：电力公用事业板块核心标的财务预测与估值]
\centering
\scriptsize
\renewcommand{\arraystretch}{1.1}
\begin{tabular*}{\textwidth}{@{\extracolsep{\fill}}l r r r r r r c@{}}
\toprule
\textbf{公司} & \textbf{市值} & \textbf{2026E 营收} & \textbf{2026E 净利} & \textbf{2026E PE} & \textbf{2027E PE} & \textbf{股息率} & \textbf{评级} \\
& (亿元) & (亿元/增速) & (亿元/增速) & (倍) & (倍) & (2026E) & \\
\midrule
长江电力 & 6,852 & 885.0 / 8\% & 345.0 / 10\% & 19.9 & 17.8 & 3.8\% & \textcolor{riskgreen}{超配} \\
国电电力 & 1,256 & 1,920.0 / 12\% & 82.5 / 25\% & 15.2 & 12.5 & 4.2\% & \textcolor{riskgreen}{超配} \\
华能国际 & 1,185 & 2,680.0 / 14\% & 78.3 / 32\% & 15.1 & 11.8 & 4.0\% & \textcolor{riskgreen}{超配} \\
华电国际 & 785 & 1,250.0 / 10\% & 52.6 / 28\% & 14.9 & 12.0 & 4.5\% & \textcolor{riskgreen}{超配} \\
中国核电 & 1,620 & 820.0 / 15\% & 128.5 / 20\% & 12.6 & 10.8 & 3.5\% & \textcolor{riskgreen}{超配} \\
中国广核 & 1,580 & 925.0 / 12\% & 112.0 / 15\% & 14.1 & 12.3 & 3.8\% & \textcolor{navy}{标配} \\
三峡能源 & 1,420 & 295.0 / 22\% & 78.5 / 25\% & 18.1 & 14.8 & 3.0\% & \textcolor{navy}{标配} \\
龙源电力 & 980 & 420.0 / 18\% & 65.2 / 22\% & 15.0 & 12.5 & 3.6\% & \textcolor{navy}{标配} \\
\bottomrule
\end{tabular*}
\vspace{0.5em}
\sourcenote{Wind，卖方一致预期，本研究团队预测；市场数据截至 2026 年 6 月 20 日}
\end{exhibitbox}

\begin{exhibitbox}[表 B-3：智驾链零部件板块核心标的财务预测与估值]
\centering
\scriptsize
\renewcommand{\arraystretch}{1.1}
\begin{tabular*}{\textwidth}{@{\extracolsep{\fill}}l r r r r r r c@{}}
\toprule
\textbf{公司} & \textbf{市值} & \textbf{2026E 营收} & \textbf{2026E 净利} & \textbf{2026E PE} & \textbf{2027E PE} & \textbf{2026E PEG} & \textbf{评级} \\
& (亿元) & (亿元/增速) & (亿元/增速) & (倍) & (倍) & (倍) & \\
\midrule
德赛西威 & 1,085 & 285.0 / 35\% & 28.5 / 45\% & 38.1 & 25.8 & 0.85 & \textcolor{riskgreen}{超配} \\
华阳集团 & 468 & 125.0 / 32\% & 12.8 / 50\% & 36.6 & 23.6 & 0.73 & \textcolor{riskgreen}{超配} \\
伯特利 & 385 & 82.5 / 30\% & 10.5 / 42\% & 36.7 & 25.8 & 0.87 & \textcolor{riskgreen}{超配} \\
中科创达 & 412 & 85.6 / 28\% & 8.2 / 35\% & 50.2 & 37.2 & 1.43 & \textcolor{navy}{标配} \\
拓普集团 & 598 & 285.0 / 25\% & 28.2 / 32\% & 21.2 & 16.2 & 0.66 & \textcolor{navy}{标配} \\
保隆科技 & 235 & 62.8 / 38\% & 6.5 / 52\% & 36.2 & 23.8 & 0.70 & \textcolor{riskgreen}{超配} \\
经纬恒润 & 210 & 55.2 / 35\% & 4.8 / 60\% & 43.8 & 27.4 & 0.73 & \textcolor{navy}{标配} \\
\bottomrule
\end{tabular*}
\vspace{0.5em}
\sourcenote{Wind，卖方一致预期，本研究团队预测；市场数据截至 2026 年 6 月 20 日}
\end{exhibitbox}

\begin{exhibitbox}[表 B-4：半导体设备材料板块核心标的财务预测与估值]
\centering
\scriptsize
\renewcommand{\arraystretch}{1.1}
\begin{tabular*}{\textwidth}{@{\extracolsep{\fill}}l r r r r r r c@{}}
\toprule
\textbf{公司} & \textbf{市值} & \textbf{2026E 营收} & \textbf{2026E 净利} & \textbf{2026E PE} & \textbf{2027E PE} & \textbf{2026E PEG} & \textbf{评级} \\
& (亿元) & (亿元/增速) & (亿元/增速) & (倍) & (倍) & (倍) & \\
\midrule
中微公司 & 1,520 & 118.0 / 32\% & 28.5 / 45\% & 53.3 & 35.2 & 1.18 & \textcolor{navy}{标配} \\
北方华创 & 2,450 & 285.0 / 35\% & 58.2 / 48\% & 42.1 & 27.8 & 0.88 & \textcolor{riskgreen}{超配} \\
拓荆科技 & 625 & 42.5 / 45\% & 9.8 / 60\% & 63.8 & 39.0 & 1.06 & \textcolor{navy}{标配} \\
华海清科 & 485 & 35.8 / 38\% & 8.5 / 52\% & 57.1 & 36.8 & 1.10 & \textcolor{navy}{标配} \\
盛美上海 & 368 & 32.5 / 30\% & 6.5 / 40\% & 56.6 & 40.2 & 1.42 & \textcolor{navy}{标配} \\
\midrule
沪硅产业 & 685 & 45.2 / 28\% & 5.2 / 85\% & 131.7 & 72.6 & 1.55 & \textcolor{navy}{标配} \\
安集科技 & 210 & 15.8 / 35\% & 3.2 / 45\% & 65.6 & 45.2 & 1.46 & \textcolor{navy}{标配} \\
鼎龙股份 & 285 & 32.5 / 22\% & 5.8 / 38\% & 49.1 & 35.6 & 1.29 & \textcolor{navy}{标配} \\
彤程新材 & 195 & 28.2 / 18\% & 4.5 / 32\% & 43.3 & 32.5 & 1.35 & \textcolor{riskamber}{关注} \\
\bottomrule
\end{tabular*}
\vspace{0.5em}
\sourcenote{Wind，卖方一致预期，本研究团队预测；市场数据截至 2026 年 6 月 20 日}
\end{exhibitbox}

\begin{exhibitbox}[表 B-5：医疗器械与创新药板块核心标的财务预测与估值]
\centering
\scriptsize
\renewcommand{\arraystretch}{1.1}
\begin{tabular*}{\textwidth}{@{\extracolsep{\fill}}l r r r r r r c@{}}
\toprule
\textbf{公司} & \textbf{市值} & \textbf{2026E 营收} & \textbf{2026E 净利} & \textbf{2026E PE} & \textbf{2027E PE} & \textbf{2026E PEG} & \textbf{评级} \\
& (亿元) & (亿元/增速) & (亿元/增速) & (倍) & (倍) & (倍) & \\
\midrule
迈瑞医疗 & 4,850 & 425.0 / 18\% & 128.5 / 22\% & 37.7 & 30.6 & 1.71 & \textcolor{riskgreen}{超配} \\
联影医疗 & 1,620 & 125.0 / 25\% & 28.5 / 32\% & 56.8 & 42.1 & 1.78 & \textcolor{navy}{标配} \\
新产业 & 512 & 52.5 / 22\% & 15.8 / 25\% & 32.4 & 25.9 & 1.30 & \textcolor{riskgreen}{超配} \\
开立医疗 & 268 & 38.5 / 28\% & 5.8 / 35\% & 46.2 & 34.2 & 1.32 & \textcolor{navy}{标配} \\
海泰新光 & 165 & 18.2 / 35\% & 3.8 / 42\% & 43.4 & 30.5 & 1.03 & \textcolor{navy}{标配} \\
\midrule
恒瑞医药 & 2,880 & 295.0 / 15\% & 62.5 / 22\% & 46.1 & 37.2 & 2.09 & \textcolor{navy}{标配} \\
百济神州 & 2,150 & 125.0 / 45\% & -8.5 / 收窄 & N/A & N/A & N/A & \textcolor{riskgreen}{超配} \\
信达生物 & 856 & 68.5 / 35\% & -5.2 / 收窄 & N/A & N/A & N/A & \textcolor{navy}{标配} \\
荣昌生物 & 320 & 22.5 / 55\% & -3.8 / 收窄 & N/A & N/A & N/A & \textcolor{navy}{标配} \\
\bottomrule
\end{tabular*}
\vspace{0.5em}
\sourcenote{Wind，卖方一致预期，本研究团队预测；市场数据截至 2026 年 6 月 20 日}
\end{exhibitbox}

\begin{exhibitbox}[表 B-6：动力电池龙头板块核心标的财务预测与估值]
\centering
\scriptsize
\renewcommand{\arraystretch}{1.1}
\begin{tabular*}{\textwidth}{@{\extracolsep{\fill}}l r r r r r r c@{}}
\toprule
\textbf{公司} & \textbf{市值} & \textbf{2026E 营收} & \textbf{2026E 净利} & \textbf{2026E PE} & \textbf{2027E PE} & \textbf{2026E PEG} & \textbf{评级} \\
& (亿元) & (亿元/增速) & (亿元/增速) & (倍) & (倍) & (倍) & \\
\midrule
宁德时代 & 9,850 & 5,280.0 / 22\% & 685.0 / 30\% & 14.4 & 11.0 & 0.48 & \textcolor{riskgreen}{超配} \\
比亚迪 & 7,250 & 9,250.0 / 28\% & 520.0 / 35\% & 13.9 & 10.2 & 0.40 & \textcolor{riskgreen}{超配} \\
亿纬锂能 & 1,580 & 825.0 / 25\% & 85.6 / 38\% & 18.5 & 13.2 & 0.49 & \textcolor{navy}{标配} \\
国轩高科 & 820 & 525.0 / 32\% & 28.5 / 55\% & 28.8 & 18.2 & 0.52 & \textcolor{navy}{标配} \\
欣旺达 & 625 & 785.0 / 28\% & 25.8 / 42\% & 24.2 & 17.0 & 0.58 & \textcolor{navy}{标配} \\
孚能科技 & 285 & 185.0 / 45\% & -5.2 / 收窄 & N/A & N/A & N/A & \textcolor{riskamber}{关注} \\
\bottomrule
\end{tabular*}
\vspace{0.5em}
\sourcenote{Wind，卖方一致预期，本研究团队预测；市场数据截至 2026 年 6 月 20 日}
\end{exhibitbox}

\vspace{0.3em}

\textbf{评级说明}：超配------预期未来 12 个月相对沪深 300 指数超额收益大于 15\%；标配------预期超额收益 5\%--15\%；关注------预期超额收益 -5\% 至 5\%，但存在向上催化可能；低配------预期超额收益低于 -5\%。评级基于本研究团队的综合判断，不构成具体投资建议。

\section{附录C：公募基金持仓数据（2026Q1）}

本附录基于 2026 年公募基金一季报数据，统计主动偏股型基金对各重点板块的持仓情况。持仓比例为板块市值占基金股票投资总市值的比例，超配/低配幅度为持仓比例相对于板块自由流通市值占 A 股比例的差值。

\begin{exhibitbox}[表 C-1：重点板块公募基金持仓分布（2026Q1）]
\centering
\scriptsize
\renewcommand{\arraystretch}{1.15}
\begin{tabular*}{\textwidth}{@{\extracolsep{\fill}}l r r r r r@{}}
\toprule
\textbf{板块} & \textbf{公募持仓比例} & \textbf{自由流通市值占比} & \textbf{超配/低配} & \textbf{环比变化 (pct)} & \textbf{拥挤度评分} \\
\midrule
AI 算力芯片 & 6.85\% & 1.85\% & +5.00\% & +0.85 & 88 \\
光模块 & 4.23\% & 0.82\% & +3.41\% & +0.62 & 92 \\
新能源（综合） & 9.52\% & 4.85\% & +4.67\% & -0.35 & 72 \\
\hspace{2em}--\quad{}储能 & 2.85\% & 0.95\% & +1.90\% & +0.45 & 82 \\
\hspace{2em}--\quad{}动力电池 & 3.25\% & 2.35\% & +0.90\% & -0.28 & 38 \\
\hspace{2em}--\quad{}光伏 & 2.68\% & 1.28\% & +1.40\% & -0.42 & 65 \\
半导体（综合） & 8.25\% & 3.25\% & +5.00\% & +0.78 & 75 \\
\hspace{2em}--\quad{}设备材料 & 2.15\% & 0.78\% & +1.37\% & +0.32 & 70 \\
\hspace{2em}--\quad{}功率半导体 & 1.28\% & 0.52\% & +0.76\% & +0.25 & 58 \\
医药生物（综合） & 7.85\% & 6.25\% & +1.60\% & -0.52 & 42 \\
\hspace{2em}--\quad{}创新药 & 1.85\% & 1.05\% & +0.80\% & +0.35 & 42 \\
\hspace{2em}--\quad{}医疗器械 & 1.62\% & 1.28\% & +0.34\% & +0.28 & 35 \\
电力公用事业 & 2.15\% & 2.85\% & -0.70\% & +0.45 & 28 \\
汽车（含智能驾驶） & 5.68\% & 3.68\% & +2.00\% & +0.65 & 60 \\
\hspace{2em}--\quad{}智驾零部件 & 1.85\% & 0.65\% & +1.20\% & +0.42 & 65 \\
国防军工 & 3.25\% & 1.55\% & +1.70\% & -0.15 & 58 \\
TMT（综合） & 15.85\% & 8.95\% & +6.90\% & +1.25 & 78 \\
\midrule
沪深300成份股 & 58.25\% & 45.65\% & +12.60\% & -0.85 & --- \\
\bottomrule
\end{tabular*}
\vspace{0.5em}
\sourcenote{Wind，基金 2026 年一季报，本研究团队测算；主动偏股型基金样本，共约 2,850 只}
\end{exhibitbox}

\begin{keyinsight}[公募持仓结构的核心特征]
从 2026Q1 公募持仓数据来看，有几个显著特征：

\vspace{0.2em}

1. \textbf{TMT 板块仍是第一大重仓方向}，整体超配 6.9 个百分点，其中 AI 算力芯片、光模块超配幅度最大，反映机构对 AI 主线的高度共识。

2. \textbf{新能源板块内部分化加剧}，储能被大幅加仓，动力电池、光伏持仓环比下降，显示机构在新能源内部进行结构调整，从制造端转向应用端。

3. \textbf{电力公用事业低配幅度收窄}，环比加仓 0.45 个百分点，但仍处于低配状态，说明机构对电力板块的认知正在转变，但尚未充分定价。

4. \textbf{医药板块整体低配且仍在减仓}，但创新药和医疗器械出现逆势加仓，反映机构在医药内部也在精选细分方向。

5. \textbf{"高景气 + 低拥挤"的电力运营、医疗器械、创新药等板块，公募持仓比例仍低于其基本面应有的水平，预期差显著}。
\end{keyinsight}

\begin{exhibitbox}[表 C-2：各板块前五大重仓股及公募持股市值]
\centering
\scriptsize
\renewcommand{\arraystretch}{1.1}
\begin{tabular*}{\textwidth}{@{\extracolsep{\fill}}l l r r@{}}
\toprule
\textbf{板块} & \textbf{重仓股} & \textbf{公募持股市值 (亿元)} & \textbf{占板块持仓比例} \\
\midrule
\multirow{5}{*}{半导体} & 北方华创 & 285.6 & 15.2\% \\
& 中微公司 & 198.3 & 10.5\% \\
& 海光信息 & 185.2 & 9.8\% \\
& 寒武纪 & 156.8 & 8.3\% \\
& 兆易创新 & 125.4 & 6.7\% \\
\cmidrule(lr){2-4}
\multirow{5}{*}{新能源} & 宁德时代 & 520.8 & 22.5\% \\
& 比亚迪 & 385.6 & 16.7\% \\
& 阳光电源 & 185.3 & 8.0\% \\
& 亿纬锂能 & 142.5 & 6.2\% \\
& 三峡能源 & 98.6 & 4.3\% \\
\cmidrule(lr){2-4}
\multirow{5}{*}{医药生物} & 恒瑞医药 & 195.3 & 12.5\% \\
& 迈瑞医疗 & 168.5 & 10.8\% \\
& 药明康德 & 145.2 & 9.3\% \\
& 爱尔眼科 & 95.6 & 6.1\% \\
& 百济神州 & 78.5 & 5.0\% \\
\cmidrule(lr){2-4}
\multirow{5}{*}{电力公用事业} & 长江电力 & 165.2 & 38.5\% \\
& 中国核电 & 68.5 & 16.0\% \\
& 华能国际 & 45.3 & 10.6\% \\
& 国电电力 & 38.2 & 8.9\% \\
& 三峡能源 & 32.5 & 7.6\% \\
\cmidrule(lr){2-4}
\multirow{5}{*}{智能驾驶} & 德赛西威 & 125.6 & 32.5\% \\
& 伯特利 & 68.5 & 17.7\% \\
& 华阳集团 & 52.3 & 13.5\% \\
& 拓普集团 & 48.6 & 12.6\% \\
& 中科创达 & 42.5 & 11.0\% \\
\bottomrule
\end{tabular*}
\vspace{0.5em}
\sourcenote{Wind，基金 2026 年一季报，本研究团队测算}
\end{exhibitbox}

\section{附录D：产业政策时间轴与关键事件}

本附录梳理 2025 年以来与本报告重点板块相关的重要产业政策与行业关键事件，按时间顺序排列。

\begin{exhibitbox}[表 D-1：重点产业政策与关键事件时间轴]
\centering
\scriptsize
\renewcommand{\arraystretch}{1.2}
\begin{tabularx}{\textwidth}{@{}l l X l@{}}
\toprule
\textbf{时间} & \textbf{发布部门} & \textbf{政策/事件内容} & \textbf{影响板块} \\
\midrule
2025.03 & 国务院 & 《"十五五"国家战略性新兴产业发展规划》发布，明确 AI、新能源、生物医药等重点方向 & 全科技成长板块 \\
2025.04 & 工信部 & 启动第四轮半导体设备国产化专项，支持 28nm 及以上成熟制程设备量产 & 半导体设备 \\
2025.05 & 国家发改委 & 出台《关于深化电价机制改革的指导意见》，放开燃煤发电上网电价浮动范围至 $\pm$30\% & 电力公用事业 \\
2025.06 & 国务院 & 印发《数字中国建设整体布局规划 2025--2030》，算力基础设施列为重点 & AI 算力、光模块 \\
2025.07 & 国家药监局 & 发布《创新医疗器械特别审查程序（2025 修订版）》，加速创新器械审批 & 医疗器械 \\
2025.08 & 工信部 & 发布《新能源汽车智能驾驶准入管理细则》，明确 L3 级自动驾驶准入标准 & 智驾零部件 \\
2025.09 & 财政部 & 提高半导体设备及零部件进口免税力度，扩大免税范围 & 半导体设备材料 \\
2025.10 & 医保局 & 公布 2025 年医保谈判结果，创新药平均降幅收窄至 38\% & 创新药 \\
2025.11 & 国务院 & 印发《加快新型电力系统建设实施方案》，明确 2030 年新能源装机占比目标 & 新能源、电力运营 \\
2025.12 & 中央经济工作会议 & 提出"以科技创新推动产业创新"，明确 2026 年战略性新兴产业重点任务 & 全科技成长板块 \\
\midrule
2026.01 & 工信部 & 发布《人形机器人创新发展指导意见》，提出 2027 年量产目标 & 人形机器人、零部件 \\
2026.02 & 国家能源局 & 公布 2025 年全国电力工业统计数据，全年发电量同比增长 6.8\% & 电力公用事业 \\
2026.02 & 上交所 & 发布《半导体行业信息披露指引》，强化晶圆厂、设备厂信披要求 & 半导体板块 \\
2026.03 & 全国两会 & 政府工作报告提出 2026 年 GDP 增长 5\% 左右，制造业投资增长 8\% 以上 & 制造业、高端装备 \\
2026.03 & 国家发改委 & 出台《推动大规模设备更新和消费品以旧换新实施方案》，医疗设备纳入重点范围 & 医疗器械 \\
2026.04 & 工信部 & 公布一季度电子信息制造业运行情况，集成电路产量同比增长 25.3\% & 半导体全产业链 \\
2026.04 & 公安部 & 发布《智能网联汽车道路测试与示范应用管理规范》，L3 上路试点扩围 & 智驾零部件 \\
2026.05 & 国家能源局 & 发布《新型电力系统发展蓝皮书》，明确电力系统转型路径 & 电力运营、特高压 \\
2026.05 & CDE & 发布《附条件批准上市药品上市后研究要求》，创新药审批政策进一步优化 & 创新药 \\
2026.06 & 中电联 & 发布 2026 年电力供需形势预测，预计全年用电量增长 7.5\%，创近年新高 & 电力公用事业 \\
2026.06 & 国务院 & 召开常务会议，研究推动集成电路产业高质量发展政策措施 & 半导体全产业链 \\
\bottomrule
\end{tabularx}
\vspace{0.5em}
\sourcenote{各部委官网，本研究团队整理}
\end{exhibitbox}

\vspace{0.3em}

\textbf{政策环境总体判断}：2025 年以来，科技产业政策支持力度持续加大，政策主线清晰。从"十五五"战略性新兴产业规划到具体行业政策，形成了自上而下的完整政策体系。其中，半导体、AI、新能源、智能驾驶、医疗器械是政策支持力度最大的五个方向，与本报告识别的高景气度板块高度吻合，构成了产业景气度的重要政策支撑。

\section{附录E：研究方法说明与数据来源}

\subsection{E.1 研究框架概述}

本报告采用"业绩验证---产业映射---政策加持---估值持仓"四维度交叉验证框架，系统性评估 A 股各板块景气度真实性和市场定价充分性，进而识别预期差显著的投资机会。

\vspace{0.3em}

\textbf{框架设计逻辑}：

\begin{enumerate}[leftmargin=2em]
\item \textbf{业绩验证维度}------景气度的"基本面锚"。通过财务数据验证景气度的真实性，排除纯概念炒作。
\item \textbf{产业映射维度}------景气度的"持续性验证"。通过产业链调研、行业数据、技术趋势判断景气度的可持续性。
\item \textbf{政策加持维度}------景气度的"外部催化剂"。政策是中国市场重要的景气度放大器，影响产业发展节奏和空间。
\item \textbf{估值持仓维度}------景气度的"定价程度检验"。通过估值水平、机构持仓、交易热度判断市场预期的充分性。
\end{enumerate}

四个维度相互印证、交叉验证，只有在多个维度上同时得到支撑的板块，才被认为是"高景气且未充分定价"的优质投资标的。

\subsection{E.2 景气度评分模型}

景气度评分采用百分制，由四个维度加权计算得出：

\begin{exhibitbox}[表 E-1：景气度评分体系及权重]
\centering
\small
\renewcommand{\arraystretch}{1.2}
\begin{tabularx}{\textwidth}{@{}l l r X@{}}
\toprule
\textbf{维度} & \textbf{细分指标} & \textbf{权重} & \textbf{计算方法} \\
\midrule
\multirow{4}{*}{业绩维度} & 营收增速 & 12\% & 最近一期单季营收同比增速，标准化处理 \\
& 净利润增速 & 12\% & 最近一期单季归母净利润同比增速，标准化处理 \\
& 毛利率变化 & 8\% & 毛利率同比变化幅度，反映盈利质量 \\
& ROE 水平 & 8\% & 滚动 ROE 绝对值，反映盈利能力 \\
\cmidrule(lr){2-4}
\multirow{3}{*}{产业维度} & 行业空间增速 & 12\% & 行业未来 3 年复合增速预期 \\
& 技术迭代速度 & 10\% & 技术进步速率与产品更新周期 \\
& 国产化率提升空间 & 8\% & 当前国产化率与目标国产化率差距 \\
\cmidrule(lr){2-4}
\multirow{2}{*}{政策维度} & 政策支持力度 & 10\% & 政策层级、支持力度、资金规模综合评分 \\
& 产业规划目标 & 5\% & 中长期产业规划对行业的拉动作用 \\
\cmidrule(lr){2-4}
\multirow{2}{*}{资金维度} & 北向资金持仓变化 & 8\% & 近一季度北向资金持仓比例变化 \\
& 公募基金持仓变化 & 7\% & 近一季度主动偏股基金持仓比例变化 \\
\midrule
\textbf{合计} & & \textbf{100\%} & \\
\bottomrule
\end{tabularx}
\vspace{0.5em}
\sourcenote{本研究团队构建}
\end{exhibitbox}

\vspace{0.3em}

\textbf{评分标准化方法}：所有原始指标经过标准化处理，映射到 0--100 分区间。对于增速类指标，采用分段线性标准化（增速越高得分越高，但边际递减）；对于变化类指标，正负方向分别处理（正向变化加分、负向变化减分）。

\subsection{E.3 拥挤度评分模型}

拥挤度评分同样采用百分制，由三个维度加权计算：

\begin{exhibitbox}[表 E-2：拥挤度评分体系及权重]
\centering
\small
\renewcommand{\arraystretch}{1.2}
\begin{tabularx}{\textwidth}{@{}l l r X@{}}
\toprule
\textbf{维度} & \textbf{细分指标} & \textbf{权重} & \textbf{计算方法} \\
\midrule
\multirow{2}{*}{公募持仓拥挤度} & 超配幅度 & 25\% & 主动偏股基金持仓比例与自由流通市值占比差值 \\
& 持仓变化率 & 15\% & 近一季度持仓比例变化幅度 \\
\cmidrule(lr){2-4}
\multirow{2}{*}{交易拥挤度} & 成交额占比分位 & 20\% & 板块成交额占全市场比例的近 5 年分位数 \\
& 换手率分位 & 10\% & 板块月均换手率的近 5 年分位数 \\
\cmidrule(lr){2-4}
\multirow{2}{*}{估值分位拥挤度} & PE（TTM）分位 & 20\% & 板块 PE（TTM）的近 5 年历史分位数 \\
& PB（LF）分位 & 10\% & 板块 PB（LF）的近 5 年历史分位数 \\
\midrule
\textbf{合计} & & \textbf{100\%} & \\
\bottomrule
\end{tabularx}
\vspace{0.5em}
\sourcenote{本研究团队构建}
\end{exhibitbox}

\vspace{0.3em}

\textbf{拥挤度等级划分}：80 分以上为高拥挤，50--80 分为中拥挤，50 分以下为低拥挤。

\subsection{E.4 "景气度---拥挤度"二维矩阵分类法}

基于景气度评分和拥挤度评分，构建二维矩阵，将板块划分为四类：

\begin{enumerate}[leftmargin=2em]
\item \textbf{明星型}：景气度 $\geq$ 70 分 且 拥挤度 $\geq$ 65 分------高景气、高拥挤，长期看好但需注意短期波动
\item \textbf{潜力型}：景气度 $\geq$ 70 分 且 拥挤度 < 65 分------高景气、低拥挤，预期差最大，是本报告核心推荐方向
\item \textbf{博弈型}：景气度 < 60 分 且 拥挤度 $\geq$ 65 分------低景气、高拥挤，主题炒作，建议回避
\item \textbf{陷阱型}：景气度 < 60 分 且 拥挤度 < 65 分------低景气、低拥挤，估值便宜但基本面乏善可陈，需警惕估值陷阱
\end{enumerate}

对于景气度处于 60--70 分区间的板块，结合景气趋势（边际变化方向）进行综合判断。

\subsection{E.5 数据来源清单}

本报告使用的数据来源包括但不限于：

\begin{exhibitbox}[表 E-3：主要数据来源清单]
\centering
\small
\renewcommand{\arraystretch}{1.15}
\begin{tabularx}{\textwidth}{@{}l X l@{}}
\toprule
\textbf{数据类别} & \textbf{具体内容} & \textbf{数据来源} \\
\midrule
\multirow{3}{*}{市场数据} & 股价、指数、成交量、换手率 & Wind 资讯 \\
& 北向资金持股、沪深港通资金流 & Wind 资讯 \\
& 行业估值（PE、PB、PS） & Wind 资讯，本研究团队计算 \\
\cmidrule(lr){2-3}
\multirow{3}{*}{财务数据} & 上市公司定期报告财务数据 & Wind 资讯，公司公告 \\
& 卖方一致预期（盈利预测） & Wind 资讯，汇总卖方研报 \\
& 业绩快报、业绩预告 & 交易所公告，Wind 资讯 \\
\cmidrule(lr){2-3}
\multirow{3}{*}{基金持仓数据} & 公募基金季报、半年报、年报 & 基金公司公告，Wind 资讯 \\
& 主动偏股基金持仓统计 & Wind 资讯，本研究团队测算 \\
& 行业配置比例与超低配分析 & 本研究团队测算 \\
\cmidrule(lr){2-3}
\multirow{4}{*}{行业数据} & 行业产量、销量、价格数据 & 国家统计局、行业协会 \\
& 半导体行业数据 & 中芯协、SEMI、TrendForce \\
& 新能源行业数据 & 中电联、国家能源局、乘联会 \\
& 医药行业数据 & 药监局、医保局、米内网 \\
\cmidrule(lr){2-3}
\multirow{3}{*}{政策法规} & 国家及部委产业政策 & 国务院、发改委、工信部、科技部等官网 \\
& 行业监管政策 & 各监管部门官网 \\
& 地方产业支持政策 & 各地方政府官网 \\
\cmidrule(lr){2-3}
\multirow{3}{*}{研究参考} & 卖方研究报告 & 各券商研究所 \\
& 上市公司调研纪要 & 机构调研公开信息 \\
& 行业白皮书、技术路线图 & 行业协会、研究机构 \\
\bottomrule
\end{tabularx}
\vspace{0.5em}
\sourcenote{本研究团队整理}
\end{exhibitbox}

\subsection{E.6 研究方法局限性说明}

本研究方法存在以下局限性，投资者在使用本报告结论时应予注意：

\begin{enumerate}[leftmargin=2em]
\item \textbf{滞后性风险}：财务数据和基金持仓数据存在 1--3 个月的滞后性，可能无法完全反映最新变化。
\item \textbf{样本偏差}：细分板块划分具有一定主观性，不同划分标准可能导致评分结果差异。
\item \textbf{量化模型局限}：评分模型无法完全捕捉定性因素（如管理层能力、突发黑天鹅事件等）。
\item \textbf{历史规律失效风险}：拥挤度与收益率的关系基于历史经验，市场结构变化可能导致历史规律失效。
\item \textbf{一致预期偏差}：卖方一致预期可能存在系统性偏差（如普遍偏乐观），影响预测准确性。
\end{enumerate}

本报告结论仅供投资参考，不构成具体投资建议。投资者应结合自身风险承受能力和投资目标，独立做出投资决策。

\section{附录F：免责声明}

\begin{disclosurebox}[免责声明]
\small

本报告由本研究团队基于公开资料独立撰写，所载信息均来源于我们认为可靠的公开渠道，但本研究团队不对这些信息的准确性和完整性作出任何保证。报告中的观点、结论和建议仅代表报告发布当日的判断，可能随时调整。

\vspace{0.3em}

本报告所载的所有内容（包括但不限于文字、数据、图表）仅供参考，不构成对任何人的投资建议、要约或承诺。投资者据此入市，风险自担。在任何情况下，本研究团队不对任何人因使用本报告任何内容所引致的任何损失承担任何责任。

\vspace{0.3em}

本报告版权归本研究团队所有。未经书面授权，任何机构和个人不得以任何形式翻版、复制、刊登、发表或引用。经授权引用刊发的，应注明出处，且不得对本报告进行任何有悖原意的引用、删节和修改。

\vspace{0.3em}

\textbf{特别提示}：
\begin{itemize}[leftmargin=2em]
\item 股市有风险，投资需谨慎。
\item 本报告不构成任何证券买卖建议，投资者应独立判断并自行承担投资风险。
\item 报告中提及的具体个股仅用于行业分析和研究示例，不构成个股推荐。
\item 过往业绩不代表未来表现，市场有风险，入市需谨慎。
\end{itemize}

\vspace{0.3em}

\begin{center}\kaishu --- 全文完 ---\end{center}

\end{disclosurebox}
